% GENERATED FILE. Edit canonical lessons, then run npm run generate:books.
% canonical-source-hash: fnv1a64:75bee6044d73d0f1
% canonical-lessons: MR-C338-raksas, MR-C338-pari, MR-C338-jadu, MR-C338-cetkin, MR-C338-sap

\chapter{Ogres and Fairies}
\label{ch:mr-ogres-and-fairies}
\hlchaptermodality{338}

\begin{chapteropening}
\textbf{By the end of this chapter:} \emph{I can tell a fairy tale with an ogre, a witch and a curse.}

Complete the last lesson of chapter 338: I can tell a fairy tale with an ogre, a witch and a curse.
\end{chapteropening}

\section[\texorpdfstring{\mr{राक्षस} (rākṣas)}{राक्षस (rākṣas)}]{\mr{राक्षस} (rākṣas) — an ogre; a demon}

\label{lesson:MR-C338-raksas}



\begin{quote}
Before the new one: say the Marathi for renowned, then the Marathi for a poet.
\end{quote}

\subsection*{You'll want to know: \mr{राक्षस}}
\textbf{\mr{राक्षस}} — \emph{rākṣas} — \textquotedblleft{}an ogre; a demon\textquotedblright{}.

\begin{quote}
\mr{गोष्टीतल्या} \mr{राक्षसाला} \mr{सात} \mr{डोकी} \mr{होती}.
\end{quote}

\begin{quote}
\emph{goṣṭītalyā rākṣasālā sāt ḍokī hotī.}
\end{quote}

\begin{quote}
The ogre in the story had seven heads.
\end{quote}

Masculine.

Plural \mr{राक्षस}.

\begin{cousinweb}
\textbf{\mr{राक्षस}} (\emph{rākṣas}) is a tatsama, taken from Sanskrit \emph{rākṣasa}, 'a demon, an evil being' (Monier-Williams).
\end{cousinweb}

\begin{grammarlens}[title={What you've built}]
One more word for telling what happened.
\end{grammarlens}

\subsection*{Your turn}
\begin{itemize}
\raggedright
  \item \emph{Say it:} \emph{rākṣas}
  \item \emph{Say it:} \emph{rākṣas}, once more
  \item \emph{Say it:} \emph{goṣṭītalyā rākṣasālā sāt ḍokī hotī.} — The ogre in the story had seven heads.
  \item \emph{Recall:} say the Marathi for renowned, then the Marathi for a poet, then say \emph{rākṣas} again
\end{itemize}

\begin{tcolorbox}[breakable,colback=teal!4,colframe=teal!35!black,title={Before you move on}]
What does \mr{राक्षस} mean? (\textquotedblleft{}An ogre; a demon\textquotedblright{}.) Say it once more.
\end{tcolorbox}
\section[\texorpdfstring{\mr{परी} (parī)}{परी (parī)}]{\mr{परी} (parī) — a fairy}

\label{lesson:MR-C338-pari}



\begin{quote}
Before the new one: say the Marathi for a poet, then the Marathi for an ogre.
\end{quote}

\subsection*{You'll want to know: \mr{परी}}
\textbf{\mr{परी}} — \emph{parī} — \textquotedblleft{}a fairy\textquotedblright{}.

\begin{quote}
\mr{परीने} \mr{मुलीला} \mr{सोन्याचे} \mr{बूट} \mr{दिले}.
\end{quote}

\begin{quote}
\emph{parīne mulīlā sonyāce būṭ dile.}
\end{quote}

\begin{quote}
The fairy gave the girl golden shoes.
\end{quote}

Feminine.

\mr{परीकथा}, a fairy tale.

\begin{cousinweb}
\textbf{\mr{परी}} (\emph{parī}) is a loan from Persian \emph{parī}, 'a fairy, a winged being'.
\end{cousinweb}

\begin{grammarlens}[title={What you've built}]
One more word for telling what happened.
\end{grammarlens}

\subsection*{Your turn}
\begin{itemize}
\raggedright
  \item \emph{Say it:} \emph{parī}
  \item \emph{Say it:} \emph{parī}, once more
  \item \emph{Say it:} \emph{parīne mulīlā sonyāce būṭ dile.} — The fairy gave the girl golden shoes.
  \item \emph{Recall:} say the Marathi for a poet, then the Marathi for an ogre, then say \emph{parī} again
\end{itemize}

\begin{tcolorbox}[breakable,colback=teal!4,colframe=teal!35!black,title={Before you move on}]
What does \mr{परी} mean? (\textquotedblleft{}A fairy\textquotedblright{}.) Say it once more.
\end{tcolorbox}
\section[\texorpdfstring{\mr{जादू} (jādū)}{जादू (jādū)}]{\mr{जादू} (jādū) — magic}

\label{lesson:MR-C338-jadu}



\begin{quote}
Before the new one: say the Marathi for an ogre, then the Marathi for a fairy.
\end{quote}

\subsection*{You'll want to know: \mr{जादू}}
\textbf{\mr{जादू}} — \emph{jādū} — \textquotedblleft{}magic\textquotedblright{}.

\begin{quote}
\mr{जादूने} \mr{भोपळ्याची} \mr{गाडी} \mr{झाली}.
\end{quote}

\begin{quote}
\emph{jādūne bhopaḷyācī gāḍī jhālī.}
\end{quote}

\begin{quote}
By magic the pumpkin turned into a carriage.
\end{quote}

Feminine.

\mr{जादू} \mr{करणे}, to cast a spell.

\begin{cousinweb}
\textbf{\mr{जादू}} (\emph{jādū}) is a loan from Persian \emph{jādū}, 'magic, sorcery'.
\end{cousinweb}

\begin{grammarlens}[title={What you've built}]
One more word for telling what happened.
\end{grammarlens}

\subsection*{Your turn}
\begin{itemize}
\raggedright
  \item \emph{Say it:} \emph{jādū}
  \item \emph{Say it:} \emph{jādū}, once more
  \item \emph{Say it:} \emph{jādūne bhopaḷyācī gāḍī jhālī.} — By magic the pumpkin turned into a carriage.
  \item \emph{Recall:} say the Marathi for an ogre, then the Marathi for a fairy, then say \emph{jādū} again
\end{itemize}

\begin{tcolorbox}[breakable,colback=teal!4,colframe=teal!35!black,title={Before you move on}]
What does \mr{जादू} mean? (\textquotedblleft{}Magic\textquotedblright{}.) Say it once more.
\end{tcolorbox}
\section[\texorpdfstring{\mr{चेटकीण} (ceṭkīṇ)}{चेटकीण (ceṭkīṇ)}]{\mr{चेटकीण} (ceṭkīṇ) — a witch}

\label{lesson:MR-C338-cetkin}



\begin{quote}
Before the new one: say the Marathi for a fairy, then the Marathi for magic.
\end{quote}

\subsection*{You'll want to know: \mr{चेटकीण}}
\textbf{\mr{चेटकीण}} — \emph{ceṭkīṇ} — \textquotedblleft{}a witch\textquotedblright{}.

\begin{quote}
\mr{जंगलातल्या} \mr{चेटकिणीने} \mr{राजकन्येला} \mr{झोपवलं}.
\end{quote}

\begin{quote}
\emph{jaṅglātalyā ceṭkiṇīne rājkanyelā jhopavlaṁ.}
\end{quote}

\begin{quote}
The witch of the forest put the princess to sleep.
\end{quote}

Feminine.

The oblique is \mr{चेटकिणी}-: \mr{चेटकिणीने}, by the witch.

\begin{cousinweb}
\textbf{\mr{चेटकीण}} (\emph{ceṭkīṇ}) is the Marathi word for a witch, with the feminine ending \emph{-īṇ}; its origin is not settled.
\end{cousinweb}

\begin{grammarlens}[title={What you've built}]
One more word for telling what happened.
\end{grammarlens}

\subsection*{Your turn}
\begin{itemize}
\raggedright
  \item \emph{Say it:} \emph{ceṭkīṇ}
  \item \emph{Say it:} \emph{ceṭkīṇ}, once more
  \item \emph{Say it:} \emph{jaṅglātalyā ceṭkiṇīne rājkanyelā jhopavlaṁ.} — The witch of the forest put the princess to sleep.
  \item \emph{Recall:} say the Marathi for a fairy, then the Marathi for magic, then say \emph{ceṭkīṇ} again
\end{itemize}

\begin{tcolorbox}[breakable,colback=teal!4,colframe=teal!35!black,title={Before you move on}]
What does \mr{चेटकीण} mean? (\textquotedblleft{}A witch\textquotedblright{}.) Say it once more.
\end{tcolorbox}
\section[\texorpdfstring{\mr{शाप} (śāp)}{शाप (śāp)}]{\mr{शाप} (śāp) — a curse}

\label{lesson:MR-C338-sap}



\begin{quote}
Before the new one: say the Marathi for magic, then the Marathi for a witch.
\end{quote}

\subsection*{You'll want to know: \mr{शाप}}
\textbf{\mr{शाप}} — \emph{śāp} — \textquotedblleft{}a curse\textquotedblright{}.

\begin{quote}
\mr{शापामुळे} \mr{राजकुमार} \mr{बेडूक} \mr{झाला}.
\end{quote}

\begin{quote}
\emph{śāpāmuḷe rājkumār beḍūk jhālā.}
\end{quote}

\begin{quote}
Because of a curse the prince became a frog.
\end{quote}

Masculine.

\mr{शाप} \mr{देणे}, to curse.

\begin{cousinweb}
\textbf{\mr{शाप}} (\emph{śāp}) is a tatsama, taken from Sanskrit \emph{śāpa}, 'a curse' (Monier-Williams).
\end{cousinweb}

\begin{grammarlens}[title={What you've built}]
One more word for telling what happened.
\end{grammarlens}

\subsection*{Your turn}
\begin{itemize}
\raggedright
  \item \emph{Say it:} \emph{śāp}
  \item \emph{Say it:} \emph{śāp}, once more
  \item \emph{Say it:} \emph{śāpāmuḷe rājkumār beḍūk jhālā.} — Because of a curse the prince became a frog.
  \item \emph{Recall:} say the Marathi for magic, then the Marathi for a witch, then say \emph{śāp} again
\end{itemize}

\begin{tcolorbox}[breakable,colback=teal!4,colframe=teal!35!black,title={Before you move on}]
What does \mr{शाप} mean? (\textquotedblleft{}A curse\textquotedblright{}.) Say it once more.
\end{tcolorbox}
